\clearpage
\fancyhead[L]{\small STELLAR OCCULTATION.}
\hypertarget{native-453-sternbedeckung}{}
\begin{center}
{\large STELLAR OCCULTATION.}

\rule{0.75\linewidth}{0.4pt}

Astronomische Nachrichten. Vol. III. No. 50. Page 22. February 1824.

\rule{0.75\linewidth}{0.4pt}
\end{center}

\begin{center}
Occultation of the Pleiades on 29 August 1820.
\end{center}

The same occultation was observed in Göttingen, Bremen, Altona, Hamburg, Berlin, Vienna and Modena.

\begin{center}
Göttingen (kindly communicated in a letter from \textsc{Gauss} [to \textsc{Bessel}]).

\begingroup\small\setlength{\tabcolsep}{4pt}
\begin{tabular}{c l c c c | c c l}
$p$ & Egress & $9^{\mathrm u}$ & $46^{\mathrm m}$ & $42^{\mathrm s}49$ & $10^{\mathrm u}$ & $5^{\mathrm m}$ & $26^{\mathrm s}78 - 1.582x$ \\
$\eta$ & --- & $9$ & $51$ & $37.37$ & $10$ & $8$ & $22.89 - 1.265x$ \\
$f$ & Ingress & $9$ & $40$ & $4.77$ & $10$ & $46$ & $56.01 + 0.219x$ \\
$h$ & --- & $9$ & $44$ & $11.60$ & $10$ & $49$ & $18.09 + 0.795x$ \\
\end{tabular}
\endgroup
\end{center}

[where $x$ denotes the correction to the difference in latitude between the Moon and the star.]

